% =============================================================
% Chapitre 2 : Etat de l'art et choix technologiques (~12 pages)
% Source : ../../CH2_Etat_de_l_art_et_positionnement.md
%          ../../DOSSIER_TECHNIQUE_MENAL.md (inventaire technologique, 28/08/2026)
% =============================================================
\chapter{État de l'art et choix technologiques}
\label{ch:etat-art}
\soustitrechapitre{Zero Trust, DevSecOps, ingénierie de la détection, enrichissement sémantique
et inventaire du socle}

Aux carences du chapitre précédent, des réponses existent : normes, outils, travaux de recherche.
Ce chapitre les expose, justifie par comparaison les choix retenus et identifie ce que l'état de
l'art \textbf{ne résout pas} --- c'est là que se situe la contribution du mémoire. Il le convertit
en un \textbf{inventaire des trente-trois briques structurantes du socle}, chacune portant son
critère décisif \emph{et} son prix.

\section{L'architecture Zero Trust}
\label{sec:zt}

\subsection{Du modèle périmétrique au Zero Trust}
\label{sec:zt-origine}

Le modèle de sécurité traditionnel repose sur une frontière : intérieur réputé sûr, extérieur
hostile, filtrage entre les deux. Trois raisons structurelles l'ont dégradé --- ressources
dispersées, utilisateurs connectés depuis n'importe où, attaques les plus coûteuses issues d'un
accès légitime détourné ---, et la question n'est plus \og comment empêcher d'entrer ? \fg{} mais
\og que peut faire un intrus une fois entré ? \fg{} Le terme \emph{Zero Trust},
popularisé par Forrester Research dès 2010 \cite{kindervag2010} et mis en œuvre à grande échelle
par Google sous le nom \emph{BeyondCorp} à partir de 2014 \cite{beyondcorp2014}, est formalisé en
2020 par le NIST (\textbf{SP~800-207} \cite{nist800207}), aujourd'hui la référence normative : la
confiance n'est jamais accordée implicitement, elle est évaluée à chaque accès selon l'identité
du demandeur, l'état de son équipement et le contexte.

\subsection{Les sept principes directeurs}
\label{sec:zt-principes}

La publication \cite{nist800207} énonce sept principes directeurs --- trois sur les ressources et
les échanges, deux sur la décision d'accès, deux sur la vérification et l'observation --- que le
\cref{tab:zt-principes} reprend avec leur traduction pour le projet : c'est la grille d'exigences
que le \cref{ch:validation} confronte à des mesures, et le premier engagement vérifiable du
mémoire.

\begingroup
\centering
\captionof{table}{Les sept principes du NIST SP 800-207 et leur traduction pour le projet}
\label{tab:zt-principes}
\scriptsize
\begin{tabular}{|P{0.85cm}|P{6.90cm}|P{7.55cm}|}
\hline
\textbf{\#} & \textbf{Principe (reformulé)} & \textbf{Traduction pour le projet} \\
\hline
P1 & Toute ressource de données ou de calcul est une ressource à protéger &
Base de données, journaux, images et secrets sont des ressources contrôlées \\ \hline
P2 & Toute communication est sécurisée, quelle que soit sa position réseau &
Le chiffrement en transit ne dépend pas d'être \og à l'intérieur \fg{} \\ \hline
P3 & L'accès est accordé par session, jamais définitivement &
Jetons à durée courte, jamais de clé permanente \\ \hline
P4 & L'accès est déterminé par une politique dynamique (identité, contexte) &
Une identité par charge de travail, droits au cas par cas \\ \hline
P5 & L'organisation mesure et surveille en continu ses ressources &
Journalisation centralisée, analyse de posture \\ \hline
P6 & Authentification et autorisation appliquées avant chaque accès &
Vérification à chaque requête, jamais une seule fois à l'entrée \\ \hline
P7 & L'organisation collecte le plus d'informations possible &
Détection, mesure et amélioration à partir des observations \\
\hline
\end{tabular}
\endgroup
\medskip

\subsection{Modèle de maturité et limites reconnues}
\label{sec:zt-limites}

Les principes disent ce qu'il faut atteindre, non où l'on en est. Le \textbf{modèle de maturité}
de la CISA \cite{cisa_zt_maturity_model} comble cet écart, et il est retenu pour une raison
unique : c'est le seul modèle public qui distingue les \emph{piliers}
fonctionnels --- identités, équipements, réseaux, applications et charges de travail, données ---
des \emph{capacités transverses} --- visibilité et analyse, automatisation et orchestration,
gouvernance ---, qui traversent tous les piliers au lieu d'en constituer un. La distinction évite
le biais le plus commun d'une auto-évaluation : compenser un pilier vide par un autre
particulièrement mûr.

\textbf{Aucun niveau n'est attribué ici} : la présente section installe la grille, l'évaluation
pilier par pilier relevant du \cref{ch:validation}. Deux précautions accompagnent son emploi :
conçu pour des agences fédérales américaines, ses niveaux ne sont pas des seuils opposables et
l'évaluation qui en découle reste \textbf{déclarative} ; et il évalue une architecture, non un
produit, de sorte qu'un pilier noté bas peut l'être par une contrainte de périmètre plutôt que
par une carence de conception.

Trois limites du modèle lui-même doivent être mentionnées : le \textbf{coût de mise en œuvre}
(composants supplémentaires, latence, points de défaillance additionnels) ; le \textbf{déplacement
du risque vers le fournisseur d'identité}, dont la compromission devient l'unique point de
défaillance majeur ; l'\textbf{écart entre discours commercial et contenu technique}, un produit
portant le nom \og Zero Trust \fg{} ne rendant pas un système Zero Trust --- une déclinaison
pratique de 2025 \cite{nist180035} confirmant qu'aucune mise en œuvre n'est valable pour tous les
contextes. Le projet retient l'approche progressive et la limite du point de défaillance unique,
que le \cref{ch:validation} éprouve par un test de compromission d'identité de service.

\section{L'approche DevSecOps}
\label{sec:devsecops}

\subsection{Principe et portes de contrôle automatisées}
\label{sec:devsecops-portes}

Le DevOps a réduit le délai entre écriture du code et mise en production, au prix d'un goulot
d'étranglement : la sécurité, vérifiée jusque-là par audit. Le DevSecOps déplace les contrôles vers
l'amont (\emph{shift-left} \cite{shiftleft_smith2001}) et les automatise. Le point souvent négligé
est le caractère \textbf{bloquant} du contrôle : un avertissement qui n'empêche pas le déploiement
ne change rien au comportement des équipes \cite{owasp_cicd_top10}.

Quatre familles de contrôles sont couramment intégrées à une chaîne de livraison, chacune posant
une question distincte. La \textbf{détection de secrets} (Gitleaks, TruffleHog) demande si un
secret a été versionné \cite{gitguardian_sprawl} ; l'\textbf{analyse statique} ou SAST (Semgrep,
CodeQL, SonarQube), si le code contient une vulnérabilité connue ; l'\textbf{analyse de
composition et d'images} ou SCA (Trivy, Grype, Dependency-Check), si un composant vulnérable est
utilisé ; l'\textbf{analyse de l'infrastructure en code} (Checkov, tfsec, Terrascan), si la
ressource sera créée de façon non sécurisée --- famille qui n'existe que parce que
l'infrastructure est décrite en code \cite{iac_morris2020}, donc analysable avant existence
(\cref{ch:realisation}). Ces quatre questions commandent le choix des quatre briques de la chaîne
de livraison (\cref{sec:inventaire}). La grille complète en compte huit, dont le
\cref{ch:realisation} dresse l'état ; les deux familles de provenance relèvent du
\cref{sec:provenance-artefacts}.

\subsection{Provenance, intégrité et référentiels}
\label{sec:provenance-artefacts}
\label{sec:referentiels}

La \textbf{provenance} des artefacts --- prouver qu'un artefact déployé provient du code et de la
chaîne attendus --- est une préoccupation récente, rendue concrète par des compromissions
documentées \cite{cisa_aa20_352a}. Trois réponses coexistent : la \textbf{nomenclature logicielle}
(SPDX, CycloneDX), la \textbf{signature d'artefacts} (Sigstore) et le cadre \textbf{SLSA}
\cite{slsa_framework} ; toutes trois ont été évaluées puis écartées (décision D07) au profit d'un
chemin moins coûteux --- étiquette liée à l'empreinte du commit et registre unique ; motif et
condition de réévaluation au \cref{sec:composants-ecartes}.

Pour les risques applicatifs, la référence est le classement de l'OWASP \cite{owasp2025} : sa
dernière édition place le contrôle d'accès défaillant en tête et la mauvaise configuration en
deuxième position, contre la cinquième en 2021 --- ce qui valide le constat du
\cref{ch:cadre-existant} ---, et ajoute une catégorie pour la chaîne d'approvisionnement
logicielle, faisant des contrôles du \cref{sec:devsecops-portes} une exigence de référentiel.
Pour la plateforme, c'est le \textbf{CIS GCP Foundation Benchmark}
\cite{cis_gcp_benchmark,gcp_container_best_practices} : plus de cent contrôles en sept domaines,
mobilisés comme \textbf{grille de configuration} et non de certification. Le mémoire publie des
écarts nommés, datés et justifiés plutôt qu'un taux de conformité, qui supposerait une grille
figée et un audit exhaustif hors périmètre. C'est ce que le projet attend d'un référentiel :
documenté et justifié, un écart est une décision d'ingénierie ; non documenté, une négligence.

\section{Supervision et ingénierie de la détection}
\label{sec:detection}

\subsection{Modèle de supervision et détection comme code}
\label{sec:supervision-modele}
\label{sec:detection-code}

Un SIEM remplit historiquement quatre fonctions : collecter des journaux hétérogènes, les
normaliser, y appliquer des règles, présenter les alertes. Son modèle économique facture le volume
ingéré, ce qui \textbf{incite à collecter moins} --- incitation qui contredit le septième principe
du Zero Trust (\cref{sec:zt-principes}) --- et son coût reste inaccessible à un budget
contraint. D'où le \emph{SIEM sur entrepôt de données} : dissocier \textbf{stockage et
interrogation} (entrepôt généraliste) de la \textbf{logique de détection} (code versionné). Il
apporte un coût aligné au volume, un langage standard, une rétention longue et aucune dépendance
à un éditeur ; il ne fournit ni règle, ni interface, ni corrélation, ni support. Un compromis,
non une supériorité.

Le format \textbf{Sigma} \cite{sigma_spec} est à la détection ce que SQL est aux bases : une
manière neutre d'exprimer une règle, indépendante du moteur, en YAML lisible (source, condition,
faux positifs, criticité, techniques associées), qui permet de traiter les règles comme du code
source. Mais \textbf{aucun moteur pris en charge ne cible l'entrepôt retenu ici} : le format a été
évalué puis écarté au titre du principe de justification (PR1), une dépendance de conversion pour
sept règles coûtant davantage qu'elle ne rapporte. Le socle conserve la propriété recherchée ---
une règle est un objet versionné, revu et testé --- et paie ce choix par l'absence de portabilité,
l'impossibilité d'importer une règle du corpus communautaire et un plafond de couverture lié au
volume traduisible à la main, mesuré au \cref{ch:validation}.

\subsection{La matrice des techniques d'attaque}
\label{sec:attck-matrice}

La matrice \textbf{MITRE ATT\&CK} \cite{mitre_attack,mitre_attack_design2018} recense les
comportements d'attaquants en tactiques (l'objectif) et techniques (le moyen) : vocabulaire commun
conception-détection, grille de couverture, cible de l'enrichissement
(\cref{sec:enrichissement}). Précaution méthodologique : la version~19 (28/04/2026) a scindé
la tactique d'évasion des défenses en deux, portant la matrice Entreprise à quinze tactiques,
222~techniques et 475~sous-techniques --- \textbf{toute mesure de couverture sans version indiquée
n'est donc pas comparable}, et la version employée est consignée en
annexe~\ref{ann:fiche-versions}. Le mémoire publie \textbf{l'inventaire nominatif des techniques
couvertes} plutôt qu'un taux : le dénominateur pertinent étant le référentiel intégral, un
pourcentage serait exact et sans valeur comparative.

\subsection{Étude comparative et solution retenue}
\label{sec:supervision-comparatif}

Cinq options ont été comparées (\cref{tab:supervision-cinq}) : produit SIEM commercial, service
SIEM cloud, pile de recherche libre auto-hébergée, solution libre intégrée, entrepôt de données
couplé à une détection en code. Les cinq sont réellement accessibles au projet, les deux premières
étant écartées par leur modèle économique et non par leur indisponibilité. La grille retient huit
des neuf critères évalués : le neuvième, l'absence de dépendance à un éditeur, laisse les trois
options restantes à égalité, et un critère qui ne discrimine pas n'a pas tranché.

\begingroup
\setlength{\tabcolsep}{3pt}
\footnotesize
\setlength{\LTpre}{4pt}
\setlength{\LTpost}{4pt}
\begin{longtable}{|P{6.6cm}|P{1.6cm}|P{1.6cm}|P{1.6cm}|P{1.6cm}|P{1.6cm}|}
\caption[Comparaison des cinq options de supervision sur huit critères]{Comparaison des cinq options de supervision sur les huit critères discriminants.
Légende : \pictoPleinAudit~favorable, \pictoPartielAudit~intermédiaire,
\pictoVideAudit~défavorable.}
\label{tab:supervision-cinq}\\
\hline
\textbf{Critère} & \textbf{SIEM comm.} & \textbf{SIEM cloud} & \textbf{Pile libre} &
\textbf{Libre intégrée} & \textbf{Entrepôt + code} \\
\hline
\endfirsthead
\hline
\textbf{Critère} & \textbf{SIEM comm.} & \textbf{SIEM cloud} & \textbf{Pile libre} &
\textbf{Libre intégrée} & \textbf{Entrepôt + code} \\
\hline
\endhead
\hline
\endfoot
Coût pour un faible volume & \pictoVideAudit & \pictoVideAudit & \pictoPartielAudit & \pictoPleinAudit & \pictoPleinAudit \\ \hline
Coût d'exploitation courante & \pictoPleinAudit & \pictoPleinAudit & \pictoVideAudit & \pictoPartielAudit & \pictoPartielAudit \\ \hline
Règles fournies au départ & \pictoPleinAudit & \pictoPleinAudit & \pictoPartielAudit & \pictoPleinAudit & \pictoVideAudit \\ \hline
Interface d'analyste fournie & \pictoPleinAudit & \pictoPleinAudit & \pictoPleinAudit & \pictoPleinAudit & \pictoVideAudit \\ \hline
Détection versionnée et testable & \pictoPartielAudit & \pictoPartielAudit & \pictoPartielAudit & \pictoPartielAudit & \pictoPleinAudit \\ \hline
Rétention longue maîtrisée & \pictoVideAudit & \pictoPartielAudit & \pictoVideAudit & \pictoPartielAudit & \pictoPleinAudit \\ \hline
Analyse exécutée sur les données stockées & \pictoVideAudit & \pictoPartielAudit & \pictoPartielAudit & \pictoVideAudit & \pictoPleinAudit \\ \hline
Maturité et support & \pictoPleinAudit & \pictoPleinAudit & \pictoPleinAudit & \pictoPartielAudit & \pictoVideAudit \\
\end{longtable}
\endgroup

La contrainte budgétaire du \cref{ch:cadre-existant} élimine les deux premières colonnes. Le
\textbf{critère décisif} entre les trois restantes est \textbf{unique} : appliquer un
rapprochement automatique aux alertes \emph{là où elles sont stockées}, sans les exporter ---
seule la dernière option le permet. Son prix se lit sur les quatre lignes perdues : ni règles, ni
interface d'analyste, ni support, et un coût d'exploitation à la charge du projet --- trois
chantiers du \cref{ch:realisation} et trois risques du \cref{ch:validation}. Le mode d'exécution
des charges a fait l'objet d'une comparaison de même forme (\cref{ann:modes-execution}).

\section{Enrichissement sémantique des alertes}
\label{sec:enrichissement}

Une règle de détection produit une alerte : horodatage, adresse, ressource, message. L'analyste
doit répondre à trois questions qu'elle ne contient pas --- \emph{quel comportement d'attaquant ?
est-ce grave ? que faire ?} ---, premier poste de charge à disparaître quand l'équipe se réduit à
une personne. L'automatiser, c'est associer chaque alerte à une technique ATT\&CK et afficher la
conduite à tenir.

\subsection{Approches disponibles}
\label{sec:enrichissement-approches}

Quatre familles permettent cette association ; le \cref{tab:enrichissement-approches} les oppose
par leur \emph{limite décisive} plutôt que par leurs mérites, car c'est la limite qui a tranché.

\begingroup
\centering
\captionof{table}{Approches d'association alerte \texorpdfstring{$\rightarrow$}{->} technique ATT\&CK}
\label{tab:enrichissement-approches}
\scriptsize
\begin{tabular}{|P{3.7cm}|P{5.0cm}|P{6.6cm}|}
\hline
Approche & Principe & Limite décisive \\
\hline
Mots-clés / regex & Table écrite à la main &
Ignore les reformulations ; maintenance manuelle \\ \hline
Similarité vectorielle (TF-IDF) & Statistique du texte &
Ignore le sens ; formulations équivalentes éloignées \\ \hline
Modèles de langue spécialisés & Proximité vectorielle = proximité de sens &
Déterministe, exécution locale ; pas d'explication en langue naturelle \\ \hline
Modèles génératifs & Le modèle rédige l'interprétation &
Non déterministe, coût par appel, injection d'instructions \\
\hline
\end{tabular}
\endgroup
\medskip

La troisième est retenue. Le dernier point de la quatrième ligne est un argument de
\textbf{sécurité} et non de performance : les journaux portent par construction du texte de tiers
non fiables, et les soumettre à un modèle génératif transmet une entrée contrôlée par l'attaquant
à un composant qui suit des instructions, là où un modèle produisant un vecteur numérique n'en
exécute aucune. C'est la seule décision de l'inventaire dont la condition de réévaluation est
\emph{aucune} (\cref{sec:composants-ecartes}).

\subsection{Modèles spécialisés et travail le plus proche}
\label{sec:smet-travail-proche}

Les modèles généralistes fondés sur l'architecture Transformer \cite{devlin_bert2019} traitent
mal le vocabulaire de la cybersécurité, d'où des modèles spécialisés comme \textbf{SecureBERT}
\cite{securebert2022} ; et rapprocher deux descriptions d'attaque exige une représentation par
\textbf{phrase}, l'architecture Sentence-BERT \cite{sbert2019}. \textbf{ATT\&CK-BERT} combine les
deux \cite{attackbert_model,smet2023} : un modèle de phrases entraîné pour que deux descriptions
d'une même action soient proches dans l'espace vectoriel, exécutable localement et sans appel
externe --- essentiel vu le coût et la nature des données traitées.

Le travail le plus proche est \textbf{SMET} (\emph{Semantic Mapping of CVE to ATT\&CK}), IFIP
DBSec 2023, étendu en 2024 \cite{smet2024} : il associe une vulnérabilité publiée à des techniques
d'attaque via ATT\&CK-BERT et une régression logistique. L'idée n'est donc pas nouvelle, et trois
différences précises séparent SMET de ce projet. L'\textbf{entrée} : une description de
vulnérabilité rédigée et structurée pour SMET, une alerte courte et bruitée ici. Le
\textbf{régime} : hors ligne sur un corpus constitué pour SMET, périodique en recette et sous
contrainte de latence et de coût ici. La \textbf{finalité} : SMET enrichit la connaissance d'une
vulnérabilité, ce projet \textbf{prototype le réordonnancement des vulnérabilités détectées à la
livraison selon les menaces observées}. Cette dernière différence est l'hypothèse d'apport du
mémoire --- \textbf{réintroduire l'observation d'exploitation dans la décision de développement}
---, objet de la question Q6 du \cref{ch:cadre-existant}.

\section{Choix de la plateforme d'hébergement}
\label{sec:choix-plateforme}

\label{sec:plateforme-criteres}
\label{sec:plateforme-decision}

La plateforme qui héberge les méthodes et les outils comparés jusqu'ici appelle la même
justification. L'application pilote tournait avant le projet sur un serveur unique loué chez un
hébergeur européen : la migration vers un fournisseur de cloud est une décision, non un préalable,
argumentée ici et enregistrée au \cref{ch:conception} sous la référence D00.

Six options ont été examinées (\cref{tab:plateforme-10-criteres}) : trois grands fournisseurs de
cloud, un fournisseur européen généraliste, un hébergeur européen à forfait comparable à
l'existant, et un hébergeur mauritanien. Les six sont réellement accessibles au projet, par le
budget comme par les compétences déclarées par l'entreprise d'accueil : comparer une option hors
de portée serait une comparaison de façade. Sur les dix critères évalués, la grille publiée en
retient sept, un critère qui ne discrimine pas n'ayant pas participé à la décision --- sont
retirées la région unique et documentée, la fédération d'identité sans secret permanent (reprise
comme brique au \cref{sec:trente-trois-briques}) et la maturité de l'infrastructure en code,
sur lesquelles les grands fournisseurs sont à égalité.

\begingroup
\centering
\captionof{table}{Comparaison des six options d'hébergement sur les sept critères discriminants}
\label{tab:plateforme-10-criteres}
{\small Légende : \pictoPleinAudit~favorable, \pictoPartielAudit~intermédiaire,
\pictoVideAudit~défavorable. \og{}Forfait\fg{} désigne un serveur loué à prix fixe chez un
hébergeur européen, comparable à l'existant ; \og{}Local\fg{} un hébergeur mauritanien.
\par\medskip}
\scriptsize
\setlength{\tabcolsep}{4pt}
\begin{tabular}{|P{0.5cm}|P{8.9cm}|c|c|c|c|c|c|}
\hline
\textbf{\#} & \textbf{Critère} & \textbf{GCP} & \textbf{AWS} & \textbf{Azure} &
\textbf{OVH} & \textbf{Forfait} & \textbf{Local} \\
\hline
1 & Palier gratuit couvrant l'entrepôt de supervision &
\pictoPleinAudit & \pictoPartielAudit & \pictoPartielAudit & \pictoVideAudit &
\pictoVideAudit & \pictoVideAudit \\ \hline
2 & Coût direct à périmètre de calcul comparable &
\pictoVideAudit & \pictoVideAudit & \pictoVideAudit & \pictoPartielAudit &
\pictoPleinAudit & \pictoPartielAudit \\ \hline
3 & Facturation nulle en l'absence de trafic &
\pictoPleinAudit & \pictoPleinAudit & \pictoPleinAudit & \pictoVideAudit &
\pictoVideAudit & \pictoVideAudit \\ \hline
4 & Proximité géographique de la Mauritanie &
\pictoVideAudit & \pictoVideAudit & \pictoVideAudit & \pictoPartielAudit &
\pictoPartielAudit & \pictoPleinAudit \\ \hline
5 & Briques de sécurité managées &
\pictoPleinAudit & \pictoPleinAudit & \pictoPleinAudit & \pictoVideAudit &
\pictoVideAudit & \pictoVideAudit \\ \hline
6 & Charge d'apprentissage pour une personne seule &
\pictoPartielAudit & \pictoVideAudit & \pictoVideAudit & \pictoPartielAudit &
\pictoPleinAudit & \pictoPleinAudit \\ \hline
7 & Réversibilité et absence d'enfermement &
\pictoVideAudit & \pictoVideAudit & \pictoVideAudit & \pictoPartielAudit &
\pictoPleinAudit & \pictoPartielAudit \\
\hline
\end{tabular}
\endgroup
\medskip

\textbf{Élimination.} La contrainte humaine du \cref{ch:cadre-existant} écarte les deux dernières
colonnes (critère~5) : une personne seule ne peut exploiter et maintenir à jour en même temps un
filtrage applicatif, un gestionnaire de clés, un annuaire d'identités par charge de travail et un
entrepôt de journaux, briques dont dépend la faisabilité des vingt exigences du
\cref{ch:besoins-menaces}. La contrainte économique écarte les modèles facturés au forfait
(critère~3) : le socle doit pouvoir ne rien coûter lorsqu'il ne sert pas. Deux réserves encadrent
la lecture : le critère~6 est une appréciation et non une mesure, et le critère~2 est apprécié à
partir des modèles de facturation --- usage contre forfait --- et non de prix relevés à une date
unique, de sorte qu'aucun montant comparatif n'est publié. Sources et dates de consultation
figurent en annexe~\ref{ann:etudes-comparatives},
\cref{ann:etudes-comparatives-plateforme}.

\textbf{Critère décisif.} Entre les trois fournisseurs restants, il est \textbf{unique} et découle
du \cref{sec:supervision-comparatif} : la recherche de similarité doit s'exécuter dans
l'entrepôt de supervision plutôt que dans une base spécialisée. Un seul des trois réunit dans un
même service le stockage, le langage de détection, la recherche vectorielle et un palier gratuit
couvrant le volume du projet ; les deux autres imposeraient un composant supplémentaire, ce que le
principe de justification (PR1) interdit. Google Cloud est retenu sur ce critère, et sur lui seul.

\textbf{Prix payé.} Trois critères sont perdus. Le \textbf{coût direct} (ligne~2) : un serveur à
forfait chez un hébergeur européen reste moins cher, à périmètre de calcul comparable. La
\textbf{réversibilité} (ligne~7) : l'architecture d'origine, faite de conteneurs et d'un fichier
de composition, était portable d'un hébergeur à l'autre ; le socle retenu ne l'est plus. La
\textbf{proximité géographique} (ligne~4) : aucun des fournisseurs examinés n'exploite de région
en Afrique de l'Ouest, et l'hébergeur mauritanien serait le seul à satisfaire ce critère comme le
seul à n'en satisfaire aucun des six autres. Conserver l'hébergeur à forfait aurait maintenu ce
coût et cette portabilité, au prix d'exploiter à la main les quatre briques ci-dessus : c'est la
contrainte d'une seule personne, non une préférence technique, qui a fermé cette voie. \textbf{Le
projet échange de la réversibilité et du coût direct contre des briques de sécurité qu'une
personne seule n'aurait pas pu exploiter} --- le \cref{ch:validation} en mesure la part
vérifiable.

\section{Inventaire technologique du socle}
\label{sec:inventaire}

La plateforme étant choisie, cette section énonce ce qui a été construit dessus. Un inventaire
technologique ne vaut que s'il résiste à deux dérives symétriques : la liste de dépendances, qui
noie la décision dans l'outillage, et la plaquette commerciale, qui énumère des avantages sans
nommer de contrepartie. \textbf{L'état décrit est celui relevé le 28/08/2026} sur la branche
principale du dépôt et l'environnement de recette, donc \textbf{postérieurement à la refonte du
plan réseau} appliquée en douze étapes du 11 au 25/08/2026 au titre de la décision D15
(\cref{ch:conception}).

\subsection{Le critère de sélection}
\label{sec:critere-selection}

Un mémoire d'ingénieur ne recense pas ses bibliothèques ; il recense ses \textbf{décisions
d'architecture matérialisées par un composant}. Un composant est une \textbf{brique structurante}
s'il satisfait \textbf{au moins deux des trois critères} suivants ; à défaut, il relève de la
sous-technologie.

\begin{enumerate}\setlength{\itemsep}{0pt}\setlength{\parskip}{0pt}
  \item \textbf{Arbitrage} : au moins une alternative sérieuse a été évaluée puis écartée, et le
  motif du rejet est écrit.
  \item \textbf{Portée architecturale} : son retrait modifierait l'\textbf{architecture} et non
  une seule implémentation --- retirer l'entrepôt de supervision change le système, retirer une
  bibliothèque de graphiques change une page.
  \item \textbf{Propriété portée} : elle porte une \textbf{propriété du système} --- sécurité,
  coût, disponibilité, reproductibilité --- que le mémoire revendique et devra démontrer.
\end{enumerate}

Le \cref{tab:critere-application} l'applique à cinq technologies du dépôt : les deux rejets y
sont plus instructifs que les deux admissions, et le cas frontière montre où passe la limite.

\begingroup
\centering
\captionof{table}{Application du critère de sélection : deux admissions, deux rejets et un cas
frontière}
\label{tab:critere-application}
\scriptsize
\setlength{\tabcolsep}{3pt}
\begin{tabular}{|P{4.7cm}|P{2.0cm}|P{8.6cm}|}
\hline
\textbf{Technologie} & \textbf{Brique ?} & \textbf{Pourquoi} \\
\hline
Entrepôt de données comme système de supervision & \textbf{Oui} &
Arbitrage documenté (D04) sur cinq options ; porte le coût \emph{et} la capacité de détection.
Trois critères sur trois \\ \hline
Fédération d'identité de la chaîne de livraison & \textbf{Oui} &
Supprime une classe de risque entière plutôt que de la gérer (D08) ; son retrait réintroduirait une
clé permanente dans le chemin de déploiement \\ \hline
Bibliothèque de graphiques du tableau de bord & \textbf{Non} &
Aucune alternative évaluée, aucune propriété du système n'en dépend, aucune couche touchée par son
remplacement. Très visible sur chaque capture, et pourtant zéro critère sur trois \\ \hline
Bibliothèque de génération de codes temporels & \textbf{Non} &
Implémentation d'un standard ouvert : \emph{le standard est la brique, pas la bibliothèque}. La
substituer ne changerait ni le protocole d'authentification ni sa surface \\ \hline
Pilote d'accès à la base de données & \textbf{Frontière} &
Sous-technologie : aucun arbitrage documenté, aucune couche modifiée par son remplacement. Mais son
caractère synchrone \emph{est} cité au \cref{ch:realisation}, car il justifie un choix de
conception de l'API \\
\hline
\end{tabular}
\endgroup
\medskip

Ce critère est ce qui empêche l'inventaire de dériver : \textbf{trente-trois composants le
franchissent, quand le dépôt en compte plusieurs centaines}.

\subsection{Les trente-trois briques structurantes}
\label{sec:trente-trois-briques}

L'inventaire se lit par couche, du socle vers la chaîne de livraison. Le
\cref{tab:synthese-couches} en donne la lecture d'ensemble : pour chacune des six couches, les
briques retenues et le critère qui a décidé de la couche entière. Le détail brique par brique ---
libellé, rôle, critère décisif et prix payé --- est publié en \cref{ann:fiche-versions}, suivi
dans le même ordre de la version retenue et du mode de verrouillage : les deux tableaux s'y lisent
ligne pour ligne. Le produit qui réalise chaque capacité y est nommé entre parenthèses à sa
première occurrence. \textbf{L'inventaire nominatif des alternatives évaluées puis écartées,
couche par couche}, figure en annexe~\ref{ann:etudes-comparatives},
\cref{ann:alternatives-couches}, avec les deux comparaisons formalisées.

\begingroup
\setlength{\tabcolsep}{2.5pt}
\renewcommand{\arraystretch}{1.0}
\begin{longtable}{|P{2.7cm}|P{6.3cm}|P{6.1cm}|}
\caption[Les trente-trois briques, par couche]{Les trente-trois briques structurantes, couche par
couche : ce que chaque couche apporte et le critère qui l'a décidée. Le détail brique par brique
figure en \cref{ann:fiche-versions}.}
\label{tab:synthese-couches}\\
\hline
\textbf{Couche} & \textbf{Briques retenues} & \textbf{Critère décisif de la couche} \\
\hline
\endfirsthead
\hline
\textbf{Couche} & \textbf{Briques retenues} & \textbf{Critère décisif de la couche} \\
\hline
\endhead
\hline
\endfoot
\textbf{N1a} --- Socle et infrastructure \emph{(huit)} &
Fournisseur de cloud public unique ; infrastructure décrite en code ; exécution de conteneurs sans
serveur ; réseau privé virtuel en mode personnalisé ; sortie réseau directe ; accès privé aux
services managés ; répartiteur de charge global ; filtrage applicatif managé au bord &
Aucun système d'exploitation à durcir ni cluster à exploiter --- décisif sous contrainte
mono-opérateur --- et un point d'entrée public unique, donc un point de contrôle unique \\ \hline
\textbf{N1b} --- Identité, secrets et chiffrement \emph{(six)} &
Gestion des identités et des accès ; fédération d'identité de la chaîne de livraison (D08) ; coffre
de secrets managé ; service de gestion de clés ; authentification par jeton signé et rôles (D10) ;
second facteur par code temporel &
Une identité par charge de travail : c'est ce qui remplace les pare-feux internes (PR2). La
fédération supprime la classe de risque de la clé permanente au lieu de la gérer \\ \hline
\textbf{N1c} --- Exécution, conteneurisation et applications \emph{(six)} &
Conteneurisation multi-étage ; registre d'images privé (D07) ; cadre applicatif à validation typée ;
cadre web à rendu serveur (D14) ; connecteur applicatif chiffré vers la base ; outil de migration de
schéma versionné &
Le contrôle d'accès devient visible dans la signature de chaque point d'entrée, donc auditable en
lecture ; et le schéma redevient reconstructible depuis le dépôt --- réponse au constat bloquant
B5 \\ \hline
\textbf{N1d} --- Données et entrepôt de supervision \emph{(cinq)} &
Base relationnelle managée en haute disponibilité (D03, D13) ; entrepôt de données généraliste
comme système de supervision (D04) ; requêtes planifiées de l'entrepôt comme ordonnanceur ;
stockage objet à clés gérées ; collecte de journaux et exports &
Un seul service réunit le stockage, le langage de détection et la recherche vectorielle, sans base
spécialisée (D06) ; et le filtrage \emph{au moment de la collecte} porte le levier de coût
principal (BNF3) \\ \hline
\textbf{N1e} --- Détection et enrichissement sémantique \emph{(quatre)} &
Règles de détection versionnées R1 à R7 ; référentiel de techniques d'attaque MITRE ATT\&CK ;
modèle de représentation de phrases spécialisé (D05) ; moteur d'inférence embarqué &
Un modèle qui produit un vecteur \textbf{n'exécute aucune instruction} : la chaîne de preuve ne
soumet jamais une entrée contrôlée par l'attaquant à un composant qui suit des instructions \\
\hline
\textbf{N1f} --- Chaîne de livraison et chaîne de provisionnement \emph{(quatre)} &
Forge logicielle et son moteur d'exécution ; détection de secrets versionnés (D09) ; analyse
statique du code ; analyse de vulnérabilité des images et de la configuration &
Trois portes bloquantes, dont \textbf{une seule couvre à la fois les images et la description
d'infrastructure} : un composant en moins, application directe de PR1 \\
\hline
\end{longtable}
\endgroup


\textbf{Une brique a quitté l'inventaire sans successeur.} La \emph{traduction d'adresses de
sortie} y figurait jusqu'à la refonte du réseau ; elle n'a pas été corrigée, elle est devenue
\textbf{sans objet} --- les interfaces de sortie directe n'ont pas d'adresse externe, l'accès
privé leur donne tout ce qu'elles doivent joindre, et le refus par défaut en sortie interdit toute
autre destination : \textbf{il ne reste aucun trafic qu'une passerelle de traduction aurait à
traiter}. La suppression a été prononcée sur mesure : le critère posé avant l'observation était
\emph{aucune connexion enregistrée par la passerelle}, la journalisation a été portée à
\texttt{ALL} du 15 au 22/08/2026, elle n'en a relevé aucune, et la suppression a suivi le
24/08/2026. La distinction n'est pas de vocabulaire : \textbf{corriger un composant mal réglé le
laisse dans l'inventaire} --- avec son coût, sa charge d'exploitation et sa surface ---
\textbf{tandis qu'établir qu'il est sans objet l'en retire}, ce que demande le principe de
justification (PR1). L'inventaire passe de trente-quatre à trente-trois briques, et la couche N1a
de neuf à huit.

\textbf{Le service de noms de domaine n'est pas une brique du socle}, et son absence est
délibérée. Les enregistrements qui font pointer les trois domaines vers le point d'entrée sont
créés à la main chez le bureau d'enregistrement, sans ressource correspondante dans la
description en code : c'est le \textbf{seul maillon de la chaîne d'exposition qui échappe au
principe de description en code} (PR4), écart enregistré sous la référence \textbf{É10}. Sa
conséquence est opérationnelle et vérifiée : une bascule de nom faite à la main a rompu une porte
de la chaîne de livraison, sans qu'aucun contrôle ne la signale.

\subsection{Les sept composants évalués puis écartés}
\label{sec:composants-ecartes}

Un inventaire qui n'énumère que ce qu'il contient laisse penser que tout ce qui n'y figure pas a
été oublié. Sept composants de l'état de l'art ont été évalués puis écartés au titre du principe
de justification (PR1), et \textbf{six le sont par arbitrage de coût}. Le \emph{périmètre de
service anti-exfiltration} : disproportionné pour un projet
unique, et sa pose exige un périmètre administratif indisponible. L'\emph{orchestrateur de
conteneurs et maillage de services} : une souplesse inutile ici, au prix d'un cluster à exploiter,
la plateforme fournissant déjà le chiffrement interne et l'identité par service. Le \emph{cache
distribué} : utile aux seules limitations de débit distribuées, mise à l'échelle hors objectif
(BNF7). Le \emph{déploiement multi-région} : hors périmètre déclaré au \cref{ch:cadre-existant},
coût doublé, perte de données déjà couverte par les sauvegardes. La \emph{signature d'artefacts
et la nomenclature logicielle} : remplacées par l'étiquette liée à l'empreinte du commit et le
registre unique. La \emph{passerelle d'authentification du fournisseur} : elle exige un annuaire
d'entreprise indisponible, le besoin étant couvert par l'autorisation applicative par rôles.

Le septième est le \textbf{modèle de langue génératif pour l'enrichissement}, \textbf{seul rejet
motivé par la sécurité} : sortie non déterministe et surface d'injection par instruction, toutes
deux rédhibitoires pour une chaîne de preuve. La \textbf{condition de réévaluation} de chacun est
enregistrée, avec le motif complet, en annexe~\ref{ann:conception-detaillee},
\cref{ann:composants-ecartes} : datable pour les six premiers, \textbf{\emph{aucune} pour le
septième}. Les deux natures ne se confondent pas --- un arbitrage de coût se rouvre avec le
contexte, un principe non --- et \textbf{un socle qui les traite de la même façon ne sait plus
laquelle de ses décisions est négociable}.

\subsection{Deux lectures transverses de l'inventaire}
\label{sec:lectures-transverses}

\paragraph{Première lecture : le fil conducteur.} Les trente-trois briques ne résultent pas de
trente-trois décisions indépendantes : \textbf{deux contraintes} du \cref{ch:cadre-existant} ---
une personne, quelques dizaines d'euros par mois --- et \textbf{un principe}, celui de
justification (PR1), agissent comme trois filtres successifs. L'\emph{humaine} écarte ce qui
demande une exploitation permanente : orchestrateur, maillage de services, filtrage, coffre de
secrets et pile de recherche auto-hébergés. L'\emph{économique} écarte les modèles au forfait et
la supervision commerciale. Le \emph{principe} écarte les doublons : base vectorielle séparée,
second registre d'images, ordonnanceur externe, analyseur de configuration distinct de l'analyseur
d'images. \textbf{Trois filtres, et il reste trente-trois briques : c'est peu pour un système qui
couvre le bord, l'identité, l'exécution, le réseau, les données, la détection, l'enrichissement,
l'observabilité et deux chaînes de livraison.}

\paragraph{Seconde lecture : ce qu'enseigne la colonne \og prix payé \fg{}.} Lue verticalement,
elle se répartit en trois natures, et la distinction n'est pas cosmétique.

\begin{itemize}\setlength{\itemsep}{0pt}\setlength{\parskip}{0pt}
  \item \textbf{Des prix de performance, tous mesurés} et reportés au \cref{ch:validation} :
  démarrage à froid, cadence plancher d'ordonnancement, délai de détection, durée de
  reconstruction de la base. Ils sont acceptés \emph{parce qu'ils sont connus} : un coût mesuré
  est gouvernable, un coût supposé ne l'est jamais.
  \item \textbf{Des prix de couverture, tous nommés} : cinq familles de règles de filtrage sur
  onze, aucune signature d'artefact, aucune nomenclature logicielle, aucune protection adaptative,
  aucun index vectoriel. Ce ne sont pas des aveux : ils délimitent \emph{ce que le socle ne
  prétend pas faire}, et c'est ce qui rend crédible ce qu'il prétend faire.
  \item \textbf{Un prix structurel, et il est unique} : la \textbf{réversibilité}. Le socle n'est
  plus portable d'un fournisseur à l'autre, et sa reconstruction ailleurs demanderait de remplacer
  le filtrage, l'entrepôt, la gestion des clés et la fédération d'identité.
\end{itemize}

Le troisième est \textbf{le seul prix que le projet ne peut pas réduire par un travail
supplémentaire} : les deux premières familles se referment avec du temps d'ingénieur --- on
calibre un démarrage, on écrit des règles --- là où la réversibilité \emph{augmente} avec chaque
brique ajoutée.

\paragraph{Un fait que ces trois natures ne classent pas.} La refonte du plan réseau (D15) a
\textbf{retiré deux briques du socle --- le connecteur d'accès sans serveur et la passerelle de
traduction d'adresses avec son routeur --- sans en ajouter aucune} : la sortie réseau directe
n'est pas un composant à exploiter, c'est un mode de rattachement d'instances existantes. Dans le
même geste, le socle a gagné un refus par défaut en sortie, un sous-réseau par locataire et huit
règles de pare-feu de plus, toutes ciblées par identité et par port : \textbf{une architecture qui
se simplifie et se durcit dans le même mouvement est assez rare pour être relevée}. Le solde
financier va dans le même sens : la suppression du connecteur retire un coût fixe permanent et
indépendant du trafic, que la tarification publique situe autour de treize~euros par mois, contre
un supplément d'ingestion de journaux de flux d'un à deux~euros. \textbf{Ce solde est un ordre de
grandeur de tarification publique et non une facture relevée} ; le vérifier suppose de comparer
les postes de facturation des mois encadrant la bascule, comparaison qui n'a pas été faite et
dont le mode opératoire est publié ici pour qu'un tiers puisse l'exécuter.

\section{Référentiels et méthodes mobilisés}
\label{sec:referentiels-methodes}

Les référentiels ne sont pas des ornements bibliographiques : ils subissent le même traitement que
les briques techniques, avec alternatives évaluées, critère décisif et limite assumée. Sept sont
mobilisés. Cinq l'ont déjà été là où ils servent --- référentiel normatif Zero Trust et modèle de
maturité au \cref{sec:zt}, classement des risques applicatifs et référentiel de configuration
de la plateforme au \cref{sec:provenance-artefacts} ---, et le sixième est le \textbf{règlement
européen sur la protection des données}, applicable au moins par le lieu du traitement, dont la
limite est que base légale, durées de conservation et droit d'effacement relèvent du responsable
de traitement, hors périmètre du socle.

Le septième --- la méthode d'analyse de risque --- appelle une justification propre, parce qu'il
structure le chapitre suivant. \textbf{STRIDE} \cite{shostack2014} applique six catégories de
menaces à chaque flux et à chaque composant : l'analyse technique est fine, mais ne traite ni les
enjeux métier ni les sources de risque. \textbf{EBIOS Risk Manager} \cite{anssi_ebios_rm2018}
déroule cinq ateliers, des valeurs métier au traitement du risque : il relie le risque technique
à l'enjeu métier, au prix d'une mise en œuvre plus lourde, et ses cinq ateliers ne sont pas tous
menés ici au même niveau de détail. \textbf{ISO/IEC~27005:2022} \cite{iso27005_2022} fournit un
cadre général, sans méthode opératoire directement applicable. \textbf{Le choix retenu combine
les deux premières}, sur un critère décisif unique : \emph{aucune des deux, seule, ne produit à la
fois le lien au métier et le détail par flux} ; la matrice
\emph{menace~$\rightarrow$~exigence~$\rightarrow$~contrôle~$\rightarrow$~test} qui en résulte
relève du \cref{ch:besoins-menaces}.


\section*{Conclusion du chapitre}
\addcontentsline{toc}{section}{Conclusion du chapitre}

Les réponses aux carences du \cref{ch:cadre-existant} existent, mais aucune ne s'applique telle
quelle : le référentiel Zero Trust suppose des moyens que le projet n'a pas ; la détection comme
code fournit un format neutre mais aucun moteur de conversion vers l'entrepôt retenu ; les modèles
de rapprochement sémantique n'ont pas été éprouvés sur des alertes courtes et bruitées.
L'inventaire des trente-trois briques a converti cet état de l'art en décisions, chacune avec son
critère décisif et son prix. Quatre \textbf{verrous} en résultent.

\begin{itemize}\setlength{\itemsep}{0pt}\setlength{\parskip}{0pt}
  \item \textbf{V1} (origine \cref{sec:zt-limites} ; traité aux ch.~\ref{ch:conception}
  et~\ref{ch:realisation}). Les architectures Zero Trust de référence supposent une équipe de
  sécurité ; aucune n'est dimensionnée pour une exploitation par une seule personne.
  \item \textbf{V2} (origine \cref{sec:detection-code} ; ch.~\ref{ch:realisation}
  et~\ref{ch:validation}). La détection comme code ne dispose pas de moteur de conversion vers
  l'entrepôt retenu : traduction manuelle, couverture plafonnée.
  \item \textbf{V3} (origine \cref{sec:smet-travail-proche} ; ch.~\ref{ch:validation}). Les
  modèles de rapprochement sémantique sont évalués sur des textes bien formés, non sur des alertes
  courtes et bruitées issues d'un système en exploitation.
  \item \textbf{V4} (origine \cref{sec:smet-travail-proche} ; ch.~\ref{ch:realisation}
  et~\ref{ch:validation}). Aucun travail examiné ne traite la boucle entre vulnérabilités
  détectées à la livraison et menaces observées en exploitation.
\end{itemize}

La contribution est d'ingénierie, un point par verrou : une \textbf{architecture Zero Trust
dimensionnée pour la contrainte réelle} d'une entreprise émergente (V1) ; une \textbf{chaîne de
détection versionnée} sur un entrepôt généraliste, limites publiées (V2) ; un \textbf{protocole
d'évaluation} face à deux références simples (V3) ; un \textbf{prototype de boucle} entre
livraison et détection (V4) --- les deux derniers restant à l'état de protocole publié, leurs
conditions de publication, énoncées avant tout relevé, n'étant pas réunies à la date de remise.
Le chapitre suivant construit la matrice qui relie chaque menace à un contrôle et au test qui en
démontre l'efficacité.

% --- Rétablissement des blancs par défaut autour des flottants : les chapitres
%     suivants appartiennent à d'autres équipes et ne doivent pas être modifiés. ---
\setlength{\textfloatsep}{20pt plus 2pt minus 4pt}
\setlength{\intextsep}{12pt plus 2pt minus 2pt}
\setlength{\floatsep}{12pt plus 2pt minus 2pt}
